\section{Related Work}

\paragraph{Differentiable Simulation} refers to computational frameworks where physical dynamics are fully differentiable with respect to system and control parameters. This paradigm has gained significant popularity in robotics and computer graphics, allowing for efficient, gradient-based optimization of policies~\citep{mora2021pods,huangplasticinelab,xuaccelerated,xianfluidlab,wangthin} and system parameters~\citep{pacnerf,xue20243d,ma2023learning,cao2024neuma,linomniphysgs} by back-propagating directly through the simulation. In general, these frameworks can be divided into two main categories. 
\emph{Learning-based approaches} use neural networks to approximate the forward dynamics~\citep{sanchez2020learning,pfaff2020learning,chendifferentiable} without explicit formulations. These methods are naturally differentiable through the networks and are typically trained on ground-truth simulation data to predict state transitions. Despite the efficiency, they require substantial training data and often have limited generalizability.
\emph{Analytical models}, on the other hand, adhere to classic simulation methods, \eg, finite element methods (FEM)~\citep{sifakis2012fem} and material point methods (MPM)~\citep{jiang2016material}, and calculate exact gradients of the underlying physical models. Some previous works, like DiffTaichi~\citep{hu2020difftaichi} and Warp~\citep{warp2022}, use automatic differentiation with explicit time-stepping schemes. Others, such as DiffSim~\citep{qiao2020scalable} and DiffCloth~\citep{li2022diffcloth}, derive analytical gradients and employ iterative solvers for implicit time integration. 

\paragraph{Deformable Linear Object Simulation} Recent advances in simulating DLOs use various solvers, yet none fully address the comprehensive requirements for robotic manipulation. These requirements encompass physical fidelity, coupling with different materials (rigid or soft), and differentiability. As summarized in~\Cref{tab:feature-comparison}, existing methods often necessitate trade-offs among these capabilities. Concretely, learning-based surrogates, \eg, Bi-LSTM~\citep{yan2020self}, GNN~\citep{wang2022offline}, DEFORM~\citep{chendifferentiable}, and Position-based Dynamics (PBD) approaches, \eg, XPBD~\citep{liu2023robotic}, SoftGym~\citep{lin2021softgym}, provide advantages like increased speed and differentiability. However, they often lack physical rigor and struggle to accurately model elastic energies and capture bending plasticity, limiting their realism in tasks involving stiff or permanently deformable materials. On the other hand, methods such as Elastica~\citep{naughton2021elastica}, C-IPC~\citep{li2021codimensional}, and IMC~\cite{tong2023fully}, which rely on physically grounded models, can effectively capture elastic behaviors but often encounter significant integration challenges. Notably, these standard implementations typically lack differentiability, rendering them unsuitable for gradient-based policy optimization.  Lastly, while differentiable solvers using Material Point Methods (MPM) or Spring-Mass systems, \eg, DaXBench~\cite{chen2023daxbench}, PhysTwin~\cite{jiang2025phystwin}, provide gradient information, they often face challenges with robust coupling or lack support for complex topological constraints. Since no single existing solver combines all above features, we developed a customized solver to address these gaps for high-fidelity DLO manipulation.

\paragraph{Deformable Linear Object Manipulation} 
represents a significant and long-standing challenge in robotics, with critical applications in industrial assembly, surgical robotics, and household tasks~\citep{saha2006motion,lee2021sample,laezza2021learning,yu2022global,lv2022dynamic,zhaole2024dexdlo,li2025gr}. Existing literature broadly categorizes DLO manipulation into two paradigms~\citep{laezza2021learning,laezza2021reform,monguzzi2023robotic}.
The first, explicit shape control, aims to deform the DLO into a precise geometric configuration. Works in this area often employ neural networks to learn the state transition of DLO dynamics, aiming to achieve a designated shape~\citep{yan2020self,wang2022offline,chendifferentiable}. In contrast, implicit shape control focuses on fulfilling high-level task conditions where the DLO's exact shape is not the primary objective. Such tasks include tying knots~\citep{saha2006motion}, routing cables~\citep{luo2024multistage}, and fixing cables into clips~\citep{li2018vision}.
Although they have achieved success in real-world applications, they are typically task-specific and difficult to generalize. Our work addresses this limitation by introducing a comprehensive simulation environment that facilitates learning versatile DLO manipulation policies.

